% ---
% id: protocol-compiler/rem:condfamily
% kind: rem
% title: "Deciding the conditional families"
% summary: >-
%   Each conditional family is decided by the compilation map, not the
%   validator, since the emitted family's conditions do not imply it; the
%   converse gives (T), (H), (P) and, with (S), (C). For a choice-weighted
%   physical family, only its own (S) is left to decide on the family itself.
% prominence: standard
% status: accepted
% uses: [def:weighted, def:exact, lem:weighted, def:compile, def:validator, thm:validator-soundness]
% mentions: [ex:condfamily, thm:choice-encoding, rem:exactness, ass:spec, ass:coin, req:physical-family-decision, def:chaincond, req:conditional-family-decision, rem:summation-condition, foundations/req:exact-emission]
% assumes: []
% refs: []
% provenance: newly-derived
% lean:
% ---

\begin{remark}\label{rem:condfamily}
A choice-weighted presentation is two declared things and one constructed one:
the weights, the conditional families $\{N_{u|t}\}_{u\in\Uset_t}$, and the
products the compiler forms from them --- the \emph{weighted family}, which is
the emitted one where the presentation is Alice's or where the receiver segment
is the identity, and is the \emph{physical} family handed to the chain where it
is Bob's at a non-trivial one. The compiler decides all three
(Definition~\ref{def:weighted},
Requirement~\ref{req:conditional-family-decision}) --- the weights by sign and
by an exact declared-order sum, each conditional family for exact
admissibility in the sense of Definition~\ref{def:exact}, and every product it
forms for representability, with the five conditions decided once more on the
weighted family wherever that family is the emitted one. This remark records
why the middle decision cannot be dropped in favour of the last, what it buys
where the weighted family is \emph{not} emitted, and what it does not disturb.

\emph{It is a decision of the compilation map and not a check of the
validator.} The conditional families are declared input: they are no slot of
the emitted declaration, their index $(u,t)$ is not an emitted index, and the
family the consumer receives is the one product family. So the decision is
taken where the presentation is read, before any product is formed
(Requirement~\ref{req:conditional-family-decision}), beside the decisions on the weights and not
among V1--V8. Nothing in Definition~\ref{def:validator} gains or loses a
conjunct by it, and Theorem~\ref{thm:validator-soundness}\,(i) and~(ii) pair the same checks
with the same target conditions whether or not the decision is taken: the
target conditions D(i)--D(viii) quantify over the emitted declaration, which
this decision does not change, and a candidate that meets them is refused at
no check because of it. What the decision affects is which
\emph{specifications} reach a candidate at all, which is the compilation map's
business and not the validator's.

\emph{The emitted family's conditions do not imply the conditional families'.}
Example~\ref{ex:condfamily} is two weights and two one-element conditional
families, neither of them a POVM, whose weighted family satisfies (T), (H),
(P), (C) and (S) and passes V1's map comparison and check V3 with every product
exactly representable.
The reason it can happen is visible in the arithmetic: (C) on the emitted
family is one identity, $\sum_{t}q_t\sum_uN_{u|t}=\idon{S}$, and the inner sums
enter it only through that one combination, so a deficit under one choice value
is paid for by a surplus under another. Deciding the outer identity therefore
says nothing about the inner ones, and \eqref{eq:condfam-bad} is what ``nothing''
looks like.

\emph{The converse does hold, in four of the five, and it is stated of the
weighted family --- whether or not that family is the emitted one.} Suppose the
conditional families of one presentation all carry one map and are each
\tref{exactly admissible} for it, the weights satisfy $q_t>0$ and
$\sum_tq_t=1$ with that sum rounding nowhere, and the representability
decision has passed on the products, so that every entry of the weighted
family \emph{is} the real product $q_tN_{u|t}$. Then, of that family: (T)
holds \emph{against that same map} --- the map the conditional families
carry, which for a physical presentation is that of the module presenting it
and not $\Smap$ --- the weighting acting blockwise
(Lemma~\ref{lem:weighted}\,(ii)); (H) holds, a real multiple of a Hermitian
block being Hermitian; (P) holds, $q_t>0$ and $N_{u|t}^{(s)}\succeq0$ giving
$q_tN_{u|t}^{(s)}\succeq0$, the product being exact; and the \emph{real} sum
over the alphabet that indexes it is the
identity on each block, which is the computation of
Lemma~\ref{lem:weighted}\,(i) read on the declared numbers --- it is available
as an identity of those numbers precisely because (S) and (C) were decided of
each conditional family and of the weights. What does \emph{not} follow is (S)
on the weighted family's own declared-order sum: that condition is on binary64
additions --- of weighted operands, accumulated across the boundaries between
choice values --- and no hypothesis on the operands controls them
(Remark~\ref{rem:summation-condition}). \emph{Only the order restricts, never
the condition}: the order the conditional family's (S) is read in is the
restriction of the order on the alphabet that indexes the weighted family
(Definition~\ref{def:weighted}), and the two conditions are independent in
both directions. Given (S), (C) follows, the computed sum then being the real
one. So (S) is the one of the five that still needs its own decision, and the
emission decision of Definition~\ref{def:weighted} is not made redundant by
this one.

\emph{And here the two cases part company.} Where the weighted family
\emph{is} the emitted one --- a presentation on Alice's side, or one on Bob's
where the receiver segment is the identity --- (H), (P), (C) and (S) are decided
on it again by check V3 and its map (T) again by check V1
(Definition~\ref{def:validator}), because those are the conditions the
consumer's admission rule reads and this module decides what it emits rather
than inferring it (\xref{foundations/req:exact-emission}); what the converse
buys there is not a check saved but the knowledge that the surviving refusals
are (S)'s and representability's (Remark~\ref{rem:exactness}). Where it is
\emph{not} --- a physical receiver family presented choice-weighted at a
non-trivial receiver segment, which lives on the physical receiver space and
is no $\Qsys$-operator family --- the converse is what the decision on it
rests on (Requirement~\ref{req:physical-family-decision}): (T), (H) and (P)
hold of it by the decision on its conditional families, (S) on its own
declared-order sum is the one condition decided of the family itself, and (C)
of it follows from that (S) together with the conditional (C), (S) and the
weight sum. The emitted
family at such a segment is a different object, the \tref[target povm]{target POVM} of the last
link, and its own five are (C3)'s and then V1's and V3's
(Definition~\ref{def:chaincond}).

\emph{The setting-dependent form has no second subject.} There the family that
is weighted is the single family $\{F_v\}_{v\in\Vset}$ the last link of
$\Chrec$ delivered, and the choice value is drawn \emph{after} the
measurement, with probability $P(f\,|\,t(v))$; there is no family of
measurements indexed by $f$, and $\{N_{v,f}\}_v$ at fixed $f$ is not one and
is not meant to be (Lemma~\ref{lem:weighted}\,(iv),
Theorem~\ref{thm:choice-encoding}\,(vii)). Read through the overlap of the two forms ---
$P(f\,|\,t)=q_f$ for every $t$, where the one conditional family is $\{F_v\}$
itself --- the decision above applies with that family as its subject, and it
was taken already: condition~(C3) of Definition~\ref{def:chaincond} decides
the exact admissibility of the family the last link delivered, per block. So
no conditional family of the setting-dependent form goes undecided and no
second decision is added there; what the compiler decides of $P$ is what
Definition~\ref{def:weighted} states of it.

\emph{What stays assumed.} A specification whose
conditional families are not measurements is refused: what is lost is a
declaration whose per-choice reading is false, and Example~\ref{ex:condfamily}
shows that such a declaration implements a probability of the choice value
other than the one the declared public-choice record states and
Assumption~\ref{ass:coin} reads. Assumption~\ref{ass:spec}
keeps what no function of the \tref{declared data} decides --- that the declared
dimensions, alphabets and tables describe the apparatus and the
declared measurements it names are the ones performed. The physical
receiver family as presented is decided on every route
(Requirement~\ref{req:physical-family-decision}); where it is presented in
choice-weighted form, its conditional families are read here like any other
and their five conditions are decided, so what is decided of the presented
family itself shrinks to condition~(S) on its own declared-order sum ---
from which its (C) then follows, by the converse above.
\end{remark}
